\SourceLeaf{235}%
also
\[
\begin{aligned}
\frac{\mathrm{d}X}{\mathrm{d}x}&=Z^3(-(1+uu)T+tuU)\\
\frac{\mathrm{d}X}{\mathrm{d}y}&=Z^3(-(1+uu)U+tuV)\\
\frac{\mathrm{d}Y}{\mathrm{d}x}&=Z^3(tuT-(1+tt)U)\\
\frac{\mathrm{d}Y}{\mathrm{d}y}&=Z^3(tuU-(1+tt)V)
\end{aligned}
\]
Setzen wir diese Werte in den vorhergehenden Ausdruck ein, so ergibt sich
\[
k=Z^6(TV-UU)(1+tt+uu)=Z^4(TV-UU)=\frac{TV-UU}{(1+tt+uu)^2}
\]
\Article{8}
Durch geeignete Wahl des Ursprungs und der Koordinatenachsen lässt sich leicht erreichen, dass für einen bestimmten Punkt \(A\) die Werte der Größen \(t,u,U\) verschwinden. Die ersten beiden Bedingungen sind bereits erfüllt, wenn man die Tangentialebene in diesem Punkt als Ebene der Koordinaten \(x,y\) wählt. Legt man außerdem den Ursprung dieser Koordinaten in den Punkt \(A\) selbst, so nimmt der Ausdruck für die Koordinate \(z\) offenbar die Form an
\[
z=\tfrac{1}{2}T^0xx+U^0xy+\tfrac{1}{2}V^0yy+\Omega
\]
wobei \(\Omega\) von höherer als zweiter Ordnung ist. Dreht man nun die Achsen \(x,y\) um einen Winkel \(M\), für den gilt
\[
\mathrm{tang}\,2M=\frac{2U^0}{T^0-V^0}
\]
so sieht man leicht, dass sich eine Gleichung der folgenden Form ergibt
\[
z=\tfrac{1}{2}Txx+\tfrac{1}{2}Vyy+\Omega
\]
Damit ist auch die dritte Bedingung erfüllt. Ist dies geschehen, so ist Folgendes klar.

I. Schneidet man die gekrümmte Fläche mit einer zu ihr normalen Ebene, die durch die Koordinatenachse \(x\) verläuft, so erhält man eine ebene Kurve, die im Punkt \(A\) den Krümmungsradius \(=\frac{1}{T}\) hat. Das positive oder negative Vorzeichen gibt an, ob die Kurve zu der Seite hin konkav oder konvex ist, in deren Richtung die Koordinaten \(z\) positiv sind.

II. Ebenso ist \(\frac{1}{V}\) der Krümmungsradius im Punkt \(A\) der ebenen Kurve, die beim Schnitt der gekrümmten Fläche mit der Ebene durch die Achsen \(y,z\) entsteht.
